\documentclass[12pt,a4paper]{report}

\usepackage[utf8]{inputenc}
\usepackage[T1]{fontenc}
\usepackage[french]{babel}
\usepackage{lmodern}
\usepackage{geometry}
\usepackage{graphicx}
\usepackage{amsmath}
\usepackage{tikz}
\usepackage{hyperref}
\usepackage{array}
\usepackage{booktabs}
\usepackage{xcolor}
\usepackage{float}
\usepackage{caption}
\usepackage{enumitem}

\usetikzlibrary{arrows.meta, positioning, shapes.geometric, calc}

\geometry{margin=2.5cm}
\hypersetup{
  colorlinks=true,
  linkcolor=blue,
  urlcolor=blue,
  citecolor=blue
}

\title{
  \textbf{Yollande : Assistante Vocale Personnalisée pour Interaction Homme-Machine}\\
  \large Rapport Technique de Fin de Projet
}
\author{Ingénieur Hermes Mbizi}
\date{\today}

\begin{document}

\maketitle

\begin{abstract}
Ce rapport présente la conception et la réalisation de Yollande, une assistante vocale
personnalisée développée dans le cadre du cours d'Interface Homme-Machine (IHM).
Le projet vise à réduire l'effet robotique des assistants classiques en améliorant la
fluidité conversationnelle, la personnalisation, la gestion de la latence et
l'automatisation du lancement de l'application sur un poste de travail.
\end{abstract}

\tableofcontents
\clearpage

\chapter{Introduction et Contexte}

Les assistants vocaux sont devenus des interfaces privilégiées pour simplifier
l'interaction entre un utilisateur et un système informatique. Cependant, leur
utilisation reste souvent freinée par un effet perçu comme robotique : réponses
trop rigides, latence trop visible, absence de personnalisation, mauvaise gestion
des interruptions et manque de continuité dans l'échange.

Dans le cadre du cours d'Interface Homme-Machine (IHM), ce projet propose une
assistante vocale nommée \textbf{Yollande}. Elle est conçue pour accompagner
l'Ingénieur Hermes Mbizi dans des tâches quotidiennes, notamment l'enregistrement
de rappels, la consultation de tâches, les annonces vocales et le lancement
automatique de l'application au démarrage du poste.

Le problème central n'est pas uniquement technique. Il concerne aussi la qualité
de l'expérience utilisateur : l'assistant doit réagir rapidement, parler de manière
naturelle, reconnaître son utilisateur, et éviter d'imposer une interface lourde.
L'enjeu IHM est donc de construire une interaction vocale qui ressemble davantage
à un dialogue avec un collègue de bureau qu'à une suite de commandes mécaniques.

\chapter{Objectifs du Projet}

Les objectifs principaux du projet sont les suivants :

\begin{itemize}[label=--]
  \item reconnaître le nom d'activation \textit{Yollande} avant de traiter une commande ;
  \item personnaliser l'échange pour l'Ingénieur Hermes Mbizi ;
  \item enregistrer, consulter, modifier et rappeler des tâches localement ;
  \item assurer que l'enregistrement des tâches et les rappels fonctionnent en mode offline ;
  \item réduire la latence entre la question de l'utilisateur et la réponse de l'assistante ;
  \item rendre les réponses vocales plus humaines grâce à des formulations variées ;
  \item permettre l'interruption de la parole de l'assistante lorsque l'utilisateur reprend la parole ;
  \item automatiser le lancement de l'application au démarrage de la machine ou via une icône ;
  \item fournir une base extensible pour intégrer ultérieurement une reconnaissance vocale totalement offline.
\end{itemize}

\chapter{Architecture Générale}

\section{Schéma d'architecture}

\begin{figure}[H]
  \centering
  \begin{tikzpicture}[
    node distance=1.2cm and 1.4cm,
    block/.style={rectangle, rounded corners, draw=black, fill=blue!8, minimum width=3.2cm, minimum height=1cm, align=center},
    local/.style={rectangle, rounded corners, draw=black, fill=green!10, minimum width=3.4cm, minimum height=1cm, align=center},
    external/.style={rectangle, rounded corners, draw=black, fill=orange!12, minimum width=3.4cm, minimum height=1cm, align=center},
    data/.style={cylinder, shape border rotate=90, draw=black, fill=gray!12, minimum height=1cm, minimum width=1.5cm, align=center},
    arrow/.style={-{Latex[length=3mm]}, thick}
  ]

    \node[block] (user) {Utilisateur\\Ingénieur Hermes Mbizi};
    \node[block, right=of user] (mic) {Capture audio\\Microphone};
    \node[block, right=of mic] (stt) {Speech-to-Text\\Reconnaissance vocale};
    \node[local, below=of stt] (wake) {Wake Word Detection\\Yollande / Yolande / Hollande};
    \node[local, right=of wake] (parser) {Parseur local\\Commandes de tâches};
    \node[external, above=of parser] (llm) {LLM optionnel\\Gemini si Internet disponible};
    \node[local, right=of parser] (router) {Routeur de commandes\\Dialogue et confirmations};
    \node[data, below=of router] (json) {Stockage JSON\\local};
    \node[local, right=of router] (scheduler) {Scheduler local\\Rappels anticipés et principaux};
    \node[block, above=of scheduler] (tts) {Text-to-Speech\\Voix Yollande};
    \node[block, right=of tts] (speaker) {Sortie vocale\\Haut-parleurs};
    \node[block, below=of scheduler] (os) {Actions système\\lancement, fichiers, dossiers};

    \draw[arrow] (user) -- (mic);
    \draw[arrow] (mic) -- (stt);
    \draw[arrow] (stt) -- (wake);
    \draw[arrow] (wake) -- (parser);
    \draw[arrow] (parser) -- (router);
    \draw[arrow] (stt) |- (llm);
    \draw[arrow] (llm) -- (router);
    \draw[arrow] (router) -- (json);
    \draw[arrow] (json) -- (scheduler);
    \draw[arrow] (scheduler) -- (tts);
    \draw[arrow] (router) -- (tts);
    \draw[arrow] (tts) -- (speaker);
    \draw[arrow] (router) -- (os);

  \end{tikzpicture}
  \caption{Architecture fonctionnelle de l'assistante vocale Yollande.}
  \label{fig:architecture-yollande}
\end{figure}

\section{Description des flux de données}

Le flux principal commence par la parole de l'utilisateur. Le microphone capture
le signal audio, puis le module de reconnaissance vocale transforme ce signal en
texte. Ce texte est ensuite filtré par le module de détection du mot d'activation.
Yollande ne traite la demande que si son nom est détecté ou si une conversation
est déjà ouverte.

Après activation, la phrase est transmise à la couche d'analyse. Les commandes
simples et critiques, comme l'ajout d'une tâche, la liste des tâches, la
suppression ou la validation, sont traitées localement par un parseur à règles.
Ce choix garantit que les fonctions de base restent disponibles hors ligne. Le
LLM externe, lorsqu'il est activé, joue uniquement un rôle complémentaire pour
les formulations plus libres.

Le routeur de commandes orchestre ensuite l'action : il crée une tâche, demande
une confirmation, lance une action système ou transmet un message vocal. Les
tâches sont persistées dans un fichier JSON local, ce qui supprime toute
dépendance à un service distant pour l'enregistrement et la consultation.

Le scheduler de rappels fonctionne en arrière-plan. Il lit les tâches locales,
compare leur échéance à l'heure courante et déclenche deux types d'alertes :
un rappel anticipé et un rappel principal. Le message est ensuite envoyé au
module de synthèse vocale pour être annoncé à l'utilisateur.

\chapter{Technologies Utilisées}

\begin{table}[H]
  \centering
  \begin{tabular}{p{4cm} p{9cm}}
    \toprule
    \textbf{Technologie} & \textbf{Rôle dans le projet} \\
    \midrule
    Python & Langage principal de l'application, choisi pour sa simplicité, son écosystème audio et sa rapidité de prototypage. \\
    Tkinter & Interface graphique légère permettant d'afficher les tâches, de tester le micro et de contrôler l'écoute continue. \\
    SpeechRecognition & Interface de reconnaissance vocale utilisée dans le prototype. \\
    pyttsx3 / SAPI Windows & Synthèse vocale locale ou native permettant de faire parler Yollande. \\
    JSON & Persistance locale des tâches, sans base de données externe. \\
    Gemini & LLM optionnel utilisé lorsque l'accès Internet est disponible. Les commandes de tâches restent traitées localement. \\
    PowerShell / Task Scheduler & Automatisation du lancement sous Windows. \\
    systemd / Launchd & Modèles de lancement automatique pour Linux et macOS. \\
    \bottomrule
  \end{tabular}
  \caption{Technologies utilisées dans le projet.}
\end{table}

\chapter{Fonctionnement Offline des Tâches et Rappels}

L'enregistrement des tâches et les rappels ont été conçus pour fonctionner sans
connexion Internet. Trois composants assurent cette autonomie :

\begin{itemize}[label=--]
  \item \textbf{Parseur local} : les commandes comme \og rappelle-moi de ... à ... \fg{},
  \og liste mes tâches \fg{} ou \og supprime la tâche \fg{} sont détectées par règles.
  \item \textbf{Stockage JSON local} : les tâches sont enregistrées dans un fichier local,
  généralement dans le dossier de données de l'utilisateur.
  \item \textbf{Scheduler local} : un thread vérifie périodiquement les échéances et déclenche
  les rappels sans service distant.
\end{itemize}

Le mode \texttt{ASSISTANT\_OFFLINE\_MODE=1} désactive explicitement l'usage de
Gemini. Dans ce mode, les fonctions essentielles de tâches restent disponibles.
La reconnaissance vocale totalement offline nécessiterait toutefois un moteur
local spécialisé, comme Vosk ou Whisper local, car le module SpeechRecognition
peut utiliser un service Web selon le backend sélectionné.

\chapter{Justification du Speech Analysis}

L'analyse de la parole est indispensable pour dépasser l'effet robotique. Dans
une interaction humaine, la communication ne se limite pas aux mots : elle
repose aussi sur les pauses, le rythme, les interruptions et les marqueurs de
discours.

\section{Détection des silences}

La détection des silences permet de déterminer quand l'utilisateur a terminé de
parler. Un seuil trop long crée une latence désagréable ; un seuil trop court
coupe la phrase avant sa fin. Le projet ajuste donc le temps de pause afin de
réduire l'attente tout en gardant une reconnaissance acceptable.

\section{Gestion des interruptions}

Un assistant naturel doit pouvoir s'arrêter lorsque l'utilisateur reprend la
parole. Le projet prévoit une interruption de la synthèse vocale en cours dès
qu'une nouvelle phrase utilisateur est détectée. Cette capacité rapproche le
système d'un dialogue réel.

\section{Marqueurs de discours humains}

Les formulations de Yollande utilisent des expressions naturelles comme
\og Au fait \fg{}, \og petite parenthèse \fg{} ou \og je vous préviens un peu
à l'avance \fg{}. Ces marqueurs masquent les micro-latences et évitent les
annonces trop mécaniques.

\chapter{Captures d'Écran et Logs}

% Remplacer les noms de fichiers ci-dessous par vos propres captures.

\begin{figure}[H]
  \centering
  % Exemple :
  % \includegraphics[width=0.85\textwidth]{images/interface_yollande.png}
  \fbox{\parbox{0.85\textwidth}{\centering Placeholder : capture de l'interface Tkinter de Yollande}}
  \caption{Interface principale de Yollande.}
  \label{fig:interface}
\end{figure}

\begin{figure}[H]
  \centering
  % Exemple :
  % \includegraphics[width=0.85\textwidth]{images/logs_rappel.png}
  \fbox{\parbox{0.85\textwidth}{\centering Placeholder : logs de création de tâche et de rappel}}
  \caption{Exemple de logs liés aux tâches et rappels.}
  \label{fig:logs}
\end{figure}

\begin{figure}[H]
  \centering
  % Exemple :
  % \includegraphics[width=0.85\textwidth]{images/raccourci_bureau.png}
  \fbox{\parbox{0.85\textwidth}{\centering Placeholder : raccourci bureau ou tâche planifiée Windows}}
  \caption{Automatisation du lancement au démarrage ou via icône.}
  \label{fig:deploiement}
\end{figure}

\chapter{Difficultés Rencontrées}

\section{Gestion de la latence}

La latence provient principalement de trois sources : l'écoute audio, la
reconnaissance vocale et la synthèse de la réponse. Le projet a réduit cette
latence en diminuant la durée de calibration du bruit ambiant, en réduisant le
seuil de silence et en lançant la synthèse vocale en arrière-plan.

\section{Faux positifs sur le mot de réveil}

Le mot \textit{Yollande} peut être transcrit différemment par le moteur vocal :
\textit{Yolande}, \textit{Hollande} ou \textit{yo lande}. Le système accepte
plusieurs variantes afin de limiter les échecs d'activation.

\section{Interruptions en temps réel}

Interrompre une voix synthétique est techniquement plus complexe que simplement
lancer une lecture audio. Le projet garde une référence au processus de synthèse
ou au moteur vocal afin de pouvoir l'arrêter lorsque l'utilisateur reprend la
parole.

\section{Déploiement sans terminal}

Une difficulté supplémentaire concerne le lancement automatique. Sous Windows,
le système utilise un script PowerShell et le Planificateur de tâches. Sous
Linux et macOS, des modèles systemd et Launchd sont proposés.

\chapter{Limites et Améliorations Possibles}

Les principales limites actuelles sont les suivantes :

\begin{itemize}[label=--]
  \item la reconnaissance vocale du prototype peut encore dépendre d'un service Web ;
  \item la détection d'interruption reste approximative et dépend du temps de reconnaissance ;
  \item l'interface graphique est fonctionnelle mais minimaliste ;
  \item la sécurité des actions système peut être renforcée par des profils utilisateur ;
  \item le moteur local de compréhension peut être enrichi pour davantage de formulations.
\end{itemize}

Les améliorations possibles incluent :

\begin{itemize}[label=--]
  \item intégrer un moteur STT offline comme Vosk ou Whisper local ;
  \item ajouter une vraie détection locale du wake word au niveau audio ;
  \item créer une interface plus moderne avec un statut audio visuel ;
  \item ajouter un historique conversationnel local ;
  \item produire un exécutable installable pour Windows ;
  \item ajouter des tests d'intégration automatisés pour les rappels.
\end{itemize}

\chapter{Conclusion}

Le projet Yollande montre qu'un assistant vocal peut être conçu comme une
interface homme-machine plus naturelle, plus personnalisée et plus autonome.
Les choix d'architecture privilégient la robustesse locale pour les fonctions
essentielles, notamment l'enregistrement des tâches et les rappels offline.
L'automatisation du lancement et la réduction de la latence participent à une
expérience plus fluide, adaptée à un usage quotidien.

\end{document}
